\section{Layers Are Rounds of Belief Propagation}
\label{sec:layers-rounds}

One transformer layer is one round of belief propagation. $L$ layers are $L$ rounds.
This is the content of Theorem~1.1 in~\cite{coppola2026transformers}, and it is the
claim that ties the previous three sections together.

A single round of BP proceeds in two steps. In the gather step, each variable node
copies its neighbors' current beliefs into scratch slots. In the update step, each
variable node computes a new belief from the gathered evidence. One transformer
layer does exactly this: attention is the gather step, the FFN is the update step,
and the strict alternation is architecturally enforced. The depth of the network is
not an engineering hyperparameter --- it is determined by the number of reasoning
hops required.

\subsection*{The Peirce Lower Bound}

No local algorithm can solve $N$-hop reasoning in fewer than $N$ rounds. This is
the Peirce lower bound, and it applies directly to belief propagation: on a tree of
diameter $d$, BP requires exactly $d$ rounds to propagate evidence from one end to
the other~\cite{pearl1988probabilistic}. A transformer with $L$ layers can perform
at most $L$ hops of reasoning. More layers are required for longer reasoning chains,
not for more parameters per se.

This reframes the question of why large models need depth. It is not that deeper
models have more capacity in a vague sense. It is that deeper models can perform
more rounds of inference --- they can follow longer chains of implication. A model
that needs to conclude $A \Rightarrow B \Rightarrow C \Rightarrow D$ needs at least
three layers, regardless of how wide each layer is.

\subsection*{Three-Phase Layer Structure}

Experiments classifying attention heads in GPT-2 small by behavior reveal a
three-phase structure across layers:

\begin{itemize}
\item \textbf{Layers 0--2 (continuous-heavy):} local feature extraction. Attention
heads attend to nearby tokens; the residual stream accumulates basic syntactic and
semantic features. This is the first round of BP: each node gathers information
from its immediate neighbors.

\item \textbf{Layer 3 (transition):} a shift in the dominant computation. Hub heads
begin to dominate. The BOS token starts accumulating a global context representation.

\item \textbf{Layers 4--11 (hub-dominated):} global context assembly. Hub heads
distribute the BOS representation to every token. This is iterative refinement of
a global prior --- each round incorporates the accumulated evidence from all previous
rounds.
\end{itemize}

The semantic-AND heads do not scale in number: 9 in GPT-2 small, 9 in GPT-2 medium.
What scales is communication infrastructure (hub heads), not logical reasoning
capacity per layer. This is consistent with the BP view: the number of reasoning
steps is set by depth, not width.

\subsection*{Residual Connections as Belief Accumulation}

The residual connection has a natural interpretation in this account. Each layer
adds its update to the accumulated belief, rather than replacing it. This is exactly
the structure of iterative BP: each round refines the beliefs from the previous
round. The residual stream after $l$ layers encodes the belief state after $l$ rounds
of message passing.

Without residual connections, each layer would compute its output from scratch and
the iterative refinement interpretation would break down. With residual connections,
the stream has a well-defined semantic content at every depth: it is the belief state
of the implicit factor graph after that many rounds of BP.

\subsection*{Convergence}

On tree-structured factor graphs, BP converges in exactly $\mathrm{diameter}(T)$
rounds and the resulting beliefs are the exact marginal posteriors. This is the
no-hallucination corollary of~\cite{coppola2026transformers}: a transformer with BP
weights, run for sufficiently many layers over a grounded tree-structured knowledge
base, computes exact Bayesian posteriors at every node.

On loopy graphs, BP may not converge, but in practice it does on the graph structures
that arise in typical knowledge bases. The \texttt{loopy} repository reports 100\%
convergence across 500 trials on five loopy graph structures, with mean KL divergence
below $0.0002$~\cite{coppola2026transformers}. The theoretical gap between trees and
loopy graphs is smaller in practice than in the worst case.